\section{Methodology}
\label{sec:method}

BLADE (\emph{Behavior Localization via Activation-Difference Estimation}) is a gradient-free, forward-only procedure that (\textit{\romannumeral1}) scores every scalar weight of a model's residual-\emph{writer} matrices by its contribution to a target behavior net of a generic-importance penalty estimated on ordinary text, (\textit{\romannumeral2}) selects a sparse effective set of layers and weights under a perplexity budget, and (\textit{\romannumeral3}) removes or reweights the behavior by multiplying \emph{only} the selected weights by a scalar factor $\alpha$. %Each ingredient adapts a well-established technique; \S\ref{sec:related} maps the borrowings and states the novel synthesis.

\subsection{Setup and Notation} Let a decoder-only transformer have $L$ layers and residual width $d$. Each layer $\ell$ writes its attention and MLP outputs into a shared \emph{residual stream}~\cite{elhage2021framework} through two linear maps --- the attention output projection and the MLP down-projection --- which we call the layer's \emph{residual writers} (the only weights BLADE edits). Within these matrices, we call the subset of weights that BLADE selects for a behavior the \emph{modulatory weights}. We use \emph{modulatory} as a functional analogy --- modulation rather than message transmission, in the sense of~\cite{sherman1998thalamus} --- and not an identified synaptic mechanism: zeroing the selected weights strongly suppresses the target behavior while measured language-modeling and zero-shot metrics change little at most operating points (\S\ref{sec:apply} and our experiments), which is consistent with a behavior-selective, gating-like role rather than a content-carrying (driving) one. This does not establish a mechanistic driver/modulator partition. The residual-writer matrices as a whole write their outputs additively to the residual stream; ``modulatory'' describes the observed effect of the selected subset for a given behavior, not of the full matrices.
\begin{equation}
  \mathcal{W}_\ell=\{\,W^{\mathrm{o}}_\ell\in\mathbb{R}^{d\times d_a},\;
                      W^{\mathrm{down}}_\ell\in\mathbb{R}^{d\times d_m}\,\}
\end{equation}
Here $d_a$ and $d_m$ are the attention- and MLP-branch widths. For a residual-writer matrix $W\in\mathbb{R}^{d\times d_{\text{in}}}$, row $i$ indexes the residual (output) coordinate and column $j$ the writer-input coordinate; $W_{ij}$ is a single scalar weight (the connection $j\to i$). A behavior is specified by a contrastive pair of prompt sets $A,B$ (e.g.\ sycophantic vs.\ non-sycophantic answers, or harmful vs.\ harmless instructions). The behavior-dependent quantities (direction and moments) are computed with forward passes only on a \texttt{train} split, while the generic-importance penalty $Q$ (\S\ref{sec:bladeg}) is estimated once per model on C4; layer/weight selection uses a \texttt{val} split (\S\ref{sec:els}). For the epistemic-uncertainty contrast (\S\ref{sec:contrasts}) we use a stricter three-way, \emph{entity-disjoint} split ($45/25/30\%$ of the entities within each template family): the direction and moments are fit on \texttt{train}, ELS selects on \texttt{select}, and every reported rate is measured on a held-out \texttt{eval} split whose entities appear in neither of the other two.

\subsection{Scoring Function}
\subsubsection{Behavior Representation}
Both quantities below are averaged over the same token span $T(x)$: the response tokens \emph{plus the end-of-turn token} for A/B behaviors, or the last prompt token for refusal and epistemic uncertainty.

\paragraph{Direction $r_\ell$.} The (example-weighted) mean of layer $\ell$'s \emph{block output} over the span, side $A$ minus side $B$, normalized:
\begin{equation}
  r_\ell=\tilde r_\ell/\lVert\tilde r_\ell\rVert,
  \qquad 
  \tilde r_\ell=\frac{1}{|A|}\sum_{x\in A}\overline{h_\ell(x)} - \frac{1}{|B|}\sum_{x\in B}\overline{h_\ell(x)},
\end{equation}
where $\overline{h_\ell(x)}=\frac{1}{|T(x)|}\sum_{t\in T(x)}h_\ell(x)_t$. This is the difference-in-means contrastive direction used for refusal by~\cite{arditi2024refusal} and, more broadly, in representation engineering and contrastive activation steering~\cite{zou2023repe,rimsky2024caa,marks2024geometry}.

\paragraph{Writer-input moments $\mu^{A},\mu^{B}$.} For each residual-writer matrix $W$ we take the
\emph{token-pooled} mean of its input activation over the span, per side:
\begin{equation}
  \mu^{A}_W=\frac{\sum_{x\in A}\sum_{t\in T(x)}a_W(x)_t}{\sum_{x\in A}|T(x)|},\qquad
  \Delta\mu_W=\mu^{A}_W-\mu^{B}_W\in\mathbb{R}^{d_{\text{in}}},
\end{equation}
where $\mu^{B}_W$ is defined analogously over $B$ and $a_W(x)_t$ is the input to $W$ at token $t$ (captured by a forward pre-hook). We note that the direction uses an example-weighted mean while the moments are token-weighted; the two coincide when all spans on a given side have equal length and differ otherwise --- answer spans are short but not necessarily equal across examples.

\subsubsection{Weight Score}
The \emph{behavioral term} of a scalar weight $W_{ij}$ (the connection $j\to i$) in writer $W$ of layer $\ell$ is the rectified product of the (unit) behavior direction, the weight, and the input mean-shift (the full BLADE score \eqref{eq:blade-g} nets a generic-importance penalty off this term):
\begin{equation}
  s_{ij}(W)=\bigl[\,(r_\ell)_i\; W_{ij}\; (\Delta\mu_W)_j\,\bigr]_+ 
  \label{eq:score}
\end{equation}
where $[\cdot]_+=\max(\cdot,0)$. The score has an exact linear meaning at fixed input statistics: summed over all weights, $\sum_{ij}(r_\ell)_i W_{ij}(\Delta\mu_W)_j=r_\ell^{\top} W\,\Delta\mu_W$ is the shift in the writer's \emph{own output} between conditions $A$ and $B$, measured along the behavior direction. For pre-norm blocks that output enters the residual stream unchanged, so this is also the writer's direct residual-stream shift; where the branch output is normalized before the residual add (Gemma-2/3) the realized residual shift is a nonlinear function of it. In neither case is it the post-edit change in behavior, which we measure empirically. The per-weight term $c_{ij}=(r_\ell)_i W_{ij}(\Delta\mu_W)_j$ is that weight's additive contribution (zeroing $W_{ij}$ changes the shift by $-c_{ij}$), and $s_{ij}=[c_{ij}]_+$ is its rectified \emph{behavioral term}. $s_{ij}$ is large when the weight routes the behavior-discriminative input shift $(\Delta\mu_W)_j$ into the residual along $r_\ell$; the rectifier keeps only weights that push \emph{toward} the behavior. Structurally this is a weight$\times$activation importance score in the family of \cite{sun2023wanda}, with the calibration-set activation norm replaced by the contrastive mean-shift $\Delta\mu$ and projected onto $r_\ell$.

\subsubsection{Generic-Importance Penalty}
\label{sec:bladeg}
The behavioral term $s_{ij}$ ranks a weight purely by how much it writes the behavior-discriminative shift into the residual; it is blind to whether the \emph{same} weight also carries generic computation. We therefore net out a data-driven \emph{generic-importance} penalty $Q_{ij}\ge 0$, estimated forward-only on generic text, that \emph{proxies} the squared perturbation zeroing $W_{ij}$ induces at one downstream normalized coordinate. Each residual writer feeds a downstream \emph{reader}: the RMSNorm that first re-reads that writer's residual contribution. For Llama-style pre-norm blocks (Llama/Qwen/Phi here) this is the attention output projection read by its layer's post-attention norm, and the MLP down-projection read by the next layer's input norm (the final norm for the last layer); architectures with branch-output norms (Gemma-2/3, whose MLP output is normalized before the residual add) place the first reader on the branch, which we treat as an approximation. Zeroing $W_{ij}$ changes the writer's output coordinate $i$ by $-W_{ij}\,a_W(x)_{t,j}$. Passing this through the reader-norm's \emph{diagonal} gain --- a fixed-denominator, diagonal approximation to the RMSNorm Jacobian that drops its rank-one normalization term --- gives an induced perturbation $\approx(\gamma_i/\mathrm{rms}_{t})\,W_{ij}\,a_W(x)_{t,j}$ at reader coordinate $i$. Here $\gamma_i$ is the reader-norm's \emph{effective} gain on coordinate $i$ (the stored parameter, or $1+$ the stored parameter for Gemma), and $\mathrm{rms}_{t}$ is the per-token RMS of the reader's input. Its expected square over generic tokens gives
\begin{equation}
  Q^{\mathrm{g1}}_{ij}=W_{ij}^2\,\gamma_i^2\;\mathbb{E}_{(x,t)\sim\mathcal{C}}\!\Big[\tfrac{a_W(x)_{t,j}^2}{\mathrm{rms}_{t}^2}\Big],
  \qquad
  Q^{\mathrm{g0}}_{ij}=W_{ij}^2\,\mathbb{E}_{(x,t)\sim\mathcal{C}}\big[a_W(x)_{t,j}^2\big],
  \label{eq:Q}
\end{equation}
where $\mathcal{C}$ is the C4 calibration set (the same corpus used for the budget $\beta$; $65$k tokens in $2048$-token windows) and $(x,t)$ ranges over its tokens. The reader-agnostic $Q^{\mathrm{g0}}$ is the \emph{squared} weight$\times$activation importance of~\cite{sun2023wanda} (their score is $|W_{ij}|\,\lVert a_j\rVert_2$), up to a corpus-wide token-count factor; the default $Q^{\mathrm{g1}}$ (the implementation's \texttt{g1scalar} estimator) refines it with the reader's diagonal gain $\gamma_i^2/\mathrm{rms}_{t}^2$, so a large-output weight is penalized only to the extent its output actually moves the downstream normalized read.

\paragraph{Regularized score (BLADE-G).} BLADE ranks weights by the behavioral term net of this penalty, in Lagrangian form:
\begin{equation}
  S_{ij}=[c_{ij}]_+\;-\;\lambda\,Q_{ij},
  \qquad
  \lambda=\frac{\operatorname{median}\{[c_{ij}]_+ : [c_{ij}]_+>0\}}{\operatorname{median}\{Q_{ij} : Q_{ij}>0\}},
  \label{eq:blade-g}
\end{equation}
with one multiplier $\lambda$ \emph{per behavior}: its numerator is the median positive behavioral term for that behavior, while $Q$ and its median denominator are computed once per model over all residual writers (behavior-independent). We \emph{deprioritize} non-behavioral weights: any weight with $S_{ij}\le 0$ is scored $-\infty$, so the ranking draws net-positive weights first and a weight enters on its behavioral merit net of generic cost. Weights with non-positive net score are reached only if $k$ exceeds the number of positive-score weights that survive the per-matrix cap (\S\ref{sec:weightsel}). We say \emph{non-positive net score} rather than ``non-behavioral'': such a weight may still have a positive behavioral term $[c_{ij}]_+$ that the generic penalty outweighs. Setting $\lambda=0$ recovers the un-penalized variant, denoted \textbf{BLADE-B}, which we keep as an ablation; the penalized score \eqref{eq:blade-g} with $Q^{\mathrm{g1}}$ is the default and is written simply ``BLADE'' below. $S_{ij}$ is therefore a forward-only, \emph{rectified} residual-writer contribution, Lagrangian-scalarized against the same weight's diagonal RMSNorm-reader perturbation proxy on generic text. It does not by itself certify out-of-distribution transfer or capability preservation, which we assess empirically. This weight-level, generic-cost-aware selection is in the spirit of set-difference safety pruning~\cite{wei2024brittleness} and calibration-based importance~\cite{sun2023wanda}, here combining a rectified behavioral contribution with a per-weight generic penalty.

\subsection{Effective-Layer Selection (ELS)}
\label{sec:els}
BLADE restricts to a small layer set $L^\star$ via two stages. Their hyperparameters are the perplexity budget $\beta=0.05$ (max.\ relative perplexity increase, \emph{calibrated on C4} and reported on held-out WikiText, following calibration-based pruning practice~\cite{sun2023wanda,frantar2023sparsegpt}), the stopping threshold $\varepsilon=0.005$ (an absolute improvement in the rate $\pi$), and two screen/test sparsities $\rho_{\text{scr}}$ (Stage 1) and $\rho_{\text{tst}}$ (Stage 2), both $0.005$ in our runs. These probe sparsities are \emph{decoupled} from the $\rho$ finally applied (\S\ref{sec:weightsel}): tying them to the final $\rho$ (re-selecting $L^\star$ per $\rho$) changes which layers are chosen --- smaller $\rho$ admits more layers --- but did not improve the resulting behavior/perplexity trade-off in our sweeps, so we keep the fixed probe. Let $\Delta\mathrm{ppl}(\cdot)$ be the relative perplexity change and $\pi(\cdot)$ the in-distribution behavior metric on \texttt{val} --- the \texttt{select} split for the epistemic contrast --- (lower $=$ more removed). $\mathrm{Prune}(S,\rho)$ scores the weights of layer set $S$ via the BLADE score \eqref{eq:blade-g} and zeroes the global top-$\rho$ fraction (\S\ref{sec:weightsel}).

\paragraph{Orientation and gating.} For A/B behaviors we first measure the model's own preference and orient side $A$ toward the \emph{model-preferred} answer (so $\pi$ is that answer's pick-rate); a behavior whose baseline pick-rate is within $0.10$ of chance is skipped as not exhibited. The gate is specific to the A/B setting, where ``not exhibited'' means \emph{at chance}: refusal and epistemic uncertainty have an absolute target of $0$ rather than a chance level, so they are never gated and are reported as measured even when the baseline rate is low.

\paragraph{Stage 1 --- candidate pool.} Keep a layer iff pruning it \emph{alone} stays within budget:
\begin{equation}
    \mathcal{P}=\{\ell:\Delta\mathrm{ppl}(\mathrm{Prune}(\{\ell\},\rho_{\text{scr}}))\le\beta\}
\end{equation}
\paragraph{Stage 2 --- best-first joint selection.} Greedily add the layer that most reduces the joint metric, subject to the budget, until no layer improves it by more than $\varepsilon$ (Algorithm~\ref{alg:els}).
\begin{algorithm}[t]
\caption{Best-first ELS (Stage 2)}
\label{alg:els}
\begin{algorithmic}[1]
\State $\mathcal{L}\gets\varnothing$;\quad $c\gets\pi(\text{unedited})$ \Comment{current metric}
\Loop
  \State $\ell^\dagger\gets\varnothing$;\quad $b\gets c$ \Comment{best-of-round}
  \For{$\ell\in\mathcal{P}\setminus\mathcal{L}$}
    \State $m\gets\pi(\mathrm{Prune}(\mathcal{L}\cup\{\ell\},\rho_{\text{tst}}))$;\quad
           $p\gets\Delta\mathrm{ppl}(\mathrm{Prune}(\mathcal{L}\cup\{\ell\},\rho_{\text{tst}}))$
    \If{$p\le\beta$ \textbf{and} $m<b$} \State $\ell^\dagger\gets\ell$;\ $b\gets m$ \EndIf
  \EndFor
  \If{$\ell^\dagger\ne\varnothing$ \textbf{and} $b<c-\varepsilon$}
     \State $\mathcal{L}\gets\mathcal{L}\cup\{\ell^\dagger\}$;\ $c\gets b$
  \Else\ \textbf{break}
  \EndIf
\EndLoop
\State \Return $L^\star\gets\mathcal{L}$
\end{algorithmic}
\end{algorithm}
Best-first is not tied to a predefined layer order (ties break by deterministic layer index) and captures the marginal benefit of a layer \emph{conditional on those already selected}; it does not recover pairs that help only jointly with zero individual marginal gain.

\paragraph{Why joint selection.} We ablate each selected layer \emph{alone} against all of $L^\star$ \emph{jointly} at the same $\rho$; this equalizes the pruned \emph{fraction} but not the pruned count, since the joint edit ranks over more matrices. The two extremes both occur in one model (Llama-3.2-3B): refusal is localized to a single layer ($L^\star=\{13\}$, refusal rate $1.00\to0.00$, so solo and joint coincide), whereas for epistemic uncertainty ($L^\star=\{14,1\}$, $\rho=0.002$) neither layer alone lowers the held-out hedge rate --- layer $14$ leaves it at its baseline $0.54$ and layer $1$ raises it to $0.67$ --- while the joint edit reduces it to $0.17$. The six A/B behaviors fall between these two extremes: in this ablation the joint edit is below every solo ablation for all of them (largest gap sycophancy, $0.73\to0.51$; smallest self-rate-highly, $0.62\to0.50$).

\paragraph{Determinism.} Stage 2 compares generation-based metrics, so floating-point non-determinism can change which layer wins a round and whether the $\varepsilon$ stopping test fires: the same model, data and hyperparameters can yield different $L^\star$ across GPUs or library versions. When an experiment compares several edits we therefore fit $L^\star$ once and \emph{pin} it across the compared runs, rather than re-selecting it per run.

\subsection{Weight selection}
\label{sec:weightsel}
Given $L^\star$, we score all weights of $\{\mathcal{W}_\ell:\ell\in L^\star\}$ with the BLADE score \eqref{eq:blade-g} (net-negative weights ranked last) and form one global ranking, subject to a per-matrix cap $c_{\max}=0.10$ (at most $\lfloor c_{\max}\lvert W\rvert\rfloor$ weights from any one matrix enter the candidate set, so no single writer is over-pruned). This matrix-capped global set-selection follows the per-matrix set construction of \cite{wei2024brittleness}. Selecting a target fraction $\rho$ of the pool zeroes the $k$ highest-scoring weights,
\begin{equation}
  k=\max\bigl(1,\ \operatorname{round}(\rho\, N)\bigr),\qquad
  N=\sum_{\ell\in L^\star}\sum_{W\in\mathcal{W}_\ell}\lvert W\rvert,
\end{equation}
Our implementation requires $\rho\le c_{\max}$, and $k$ is additionally capped by the pooled per-matrix candidate count $\sum_W\lfloor c_{\max}\lvert W\rvert\rfloor$ (the caps use $\lfloor\cdot\rfloor$ while $k$ uses rounding, so the cap can bind). We sweep $\rho$ over a small grid ($\{0.0005,0.002,0.005,0.02\}$). For cross-model comparison we report the operating point $\rho=0.005$ for every behavior except sycophancy and corrigibility, which are instead reported at the most-removal feasible point of the swept grid, $\arg\min_{\rho}\{\pi_{\text{val}}:\Delta\mathrm{ppl}_{\mathrm{C4}}\le\beta\}$; since $\rho=0.005$ is itself feasible, that choice cannot be worse on the selection metric it minimizes, a guarantee that does not extend to held-out behavior. A common $\rho$ does \emph{not} equalize the number of edited weights across cells: $\rho$ is a fraction of $N$, which varies with the model and with $\lvert L^\star\rvert$. The budget is enforced on the C4 calibration corpus, so an individual reported cell can still exceed $\beta$ on held-out WikiText; we report the held-out number as measured rather than re-tuning the operating point to it.

\subsection{Application: reweighting by $\alpha$}
\label{sec:apply}
Let $\mathcal{S}$ be the selected weights. BLADE multiplies them by a scalar: 
\begin{equation}
  W^{\prime}_{ij}=\alpha\,W_{ij}
\end{equation}
for $(i,j)\in\mathcal{S}$, all other unselected weights unchanged. $\alpha=0$ (the main setting, identical to zeroing / pruning $\mathcal{S}$) is intended to remove the behavior; $0<\alpha<1$ to attenuate it; $\alpha=1$ is the original model; $\alpha>1$ to amplify it. Rescaling a localized, direction-aligned component by a scalar coefficient to strengthen or weaken a behavior is in the spirit of task arithmetic~\cite{ilharco2023task} and the scale coefficient of activation steering~\cite{rimsky2024caa}, here applied at the weight level. Empirically, at $\alpha=0$ and small $\alpha$ the behavior changes while the measured capability metrics (WikiText perplexity; zero-shot accuracy where measured) change little at most reported operating points. This is consistent with a behavior-selective, capability-sparing role over the evaluated $\alpha$ range, but not uniformly so: a minority of reported cells exceed the $5\%$ budget on held-out WikiText, up to $+11.5\%$. This evidence does not identify a mechanistic driver/modulator partition. Large $\lvert\alpha\rvert$ (and negative $\alpha$) can raise perplexity and fall outside the zeroing budget. The $\alpha$-sweep is a follow-up analysis; where the mask is the removal mask, $\mathcal{S}$ is budgeted for $\alpha=0$ and the budget is not re-checked at the applied $\alpha$.

\paragraph{Scale-aware selection.} The penalty admits a scale-aware form when the intended edit is not removal, in which the mask is selected under the penalty at the intended scale. Multiplying $W_{ij}$ by $\alpha$ changes its behavioral contribution from $c_{ij}$ to $\alpha c_{ij}$ (a change of $(\alpha-1)c_{ij}$) and scales its generic reader perturbation \eqref{eq:Q} by $(\alpha-1)^2$ (a squared perturbation). For \emph{amplification} ($\alpha>1$) write $\delta=\alpha-1>0$; the amplify objective $S^{\mathrm{amp}}_{ij}=\delta[c_{ij}]_+-\lambda\delta^2 Q_{ij}$ is a positive multiple of $[c_{ij}]_+-\lambda\delta Q_{ij}$, so it ranks weights by \eqref{eq:blade-g} with an effective penalty $\lambda_{\mathrm{eff}}=\lambda\delta$ that grows with the amplification strength. For \emph{attenuation/removal} ($\alpha<1$) the benefit is behavior \emph{reduction}: with $\delta'=1-\alpha>0$ the objective $\delta'[c_{ij}]_+-\lambda\delta'^2 Q_{ij}$ ranks by $[c_{ij}]_+-\lambda\delta' Q_{ij}$, i.e.\ $\lambda_{\mathrm{eff}}=\lambda\delta'$; full removal $\alpha=0$ gives $\delta'=1$ and recovers \eqref{eq:blade-g}. In both directions $\lambda_{\mathrm{eff}}=\lambda\lvert\alpha-1\rvert$. We use this scale-aware mask in the amplification studies that select a mask specifically for $\alpha>1$. The $\alpha$-sweeps that reuse a removal-selected mask --- applying raw $\alpha W$ to the $\alpha=0$ selection --- are the \emph{fixed-mask} baseline, in which the mask is not tuned to the scale at which it is applied. Selecting at the intended scale re-ranks the mask but does not by itself guarantee the budget at that scale, which we check empirically; we label which of the two protocols is used wherever an $\alpha>1$ result is reported.

\paragraph{Contrast constructions.}
\label{sec:contrasts}
Three instantiations of $(A,B)$ are used, all feeding the identical score \eqref{eq:blade-g}, ELS, and reweighting. \emph{(i) Multiple-choice (A/B):} $A/B$ are the two answers of Anthropic-style evals; the span is the answer tokens ($+$EOT) and $\pi$ is the answer pick-rate (chance $=0.5$). \emph{(ii) Refusal:} $A/B$ are harmful (AdvBench) vs.\ harmless (Alpaca) instructions; the span is the last prompt token and $\pi$ is the keyword refusal rate on generations (target $=0$). \emph{(iii) Epistemic uncertainty:} $A/B$ are questions the model cannot answer (non-existent or unknowable referents) vs.\ matched questions it can, drawn from a templated set of $306$ items ($156$ uncertain, $150$ certain) over five families (capitals, person attributes, event years, quantities, authorship); as for refusal the span is the last prompt token, and $\pi$ is the fraction of generations matching a fixed list of hedging/abstention markers (target $=0$ for removal), the same style of keyword metric used for refusal. Cases (ii) and (iii) share a form --- a last-prompt-token contrast whose metric is a generation-level rate with an absolute target --- and differ from (i), whose metric is a forced-choice pick-rate against chance. An earlier sycophancy-only variant instead built moments from biased-vs-neutral \emph{prompt} activations at the last prompt token; we do not report results from it, and every number in this paper uses the answer-span pipeline above.